\subsection{一个关于测度分解的结果}

命题 13. --- 设 X 是拓扑空间，\(\nu\) 是 X 上的适度测度，p 是从 X 到拓扑空间 T 的 \(\nu\) -逆紧映射，且 \(\mu = p(\nu)\)。假设 X 的每个紧子空间都是可度量的。则存在映射 \(t \mapsto \lambda_{t}\)，从 T 到 \(\mathcal{M}_{+}(X)\)，具有以下性质：

\begin{enumerate}
\def\labelenumi{\alph{enumi})}
\item
  对于每个 \(t \in \mathbf{T}\)，测度 \(\lambda_t\) 集中于 \(\overline{p}^{-1}(t)\)；
\item
  对于 X 上每个普遍可测的 \(^{(1)}\) 正函数 f，函数 \(t \mapsto \lambda_{t}^{\bullet}(f)\) 在 T 上普遍可测，并且
\end{enumerate}

\[
\int_ {\mathrm{X}} ^ {\bullet} f (x) d \nu (x) = \int_ {\mathrm{T}} ^ {\bullet} d \mu (t) \int_ {\mathrm{X}} ^ {\bullet} f (x) d \lambda_ {t} (x);\tag{12}
\]

\begin{enumerate}
\def\labelenumi{\alph{enumi})}
\setcounter{enumi}{2}
\tightlist
\item
  对 \(t \in \mathrm{T}\) 中满足 \(\lambda_t(1) \neq 1\) 的点所成的集合局部 \(\mu\) -可忽略。
\end{enumerate}

此外，如果 \(t \mapsto \lambda_t'\) 是从 \(\mathbf{T}\) 到 \(\mathcal{M}_+(\mathbf{X})\) 的映射，满足条件 \(a)\) 和 \(b)\)，则 \(t \in \mathbf{T}\) 中满足 \(\lambda_t \neq \lambda_t'\) 的点所成的集合局部 \(\mu\) -可忽略。我们需要一个辅助结果：

引理 4. --- 设 X 是拓扑空间，\(\nu\) 是 X 上的测度，\(f\) 是从 X 到拓扑空间 F（不论是否 Hausdorff）的 \(\nu\) -可测映射。则存在从 X 到 F 的普遍可测映射 \(f'\)，它与 \(f\) 局部 \(\nu\) -几乎处处相等。

证明与 Ch. V, §3, No.~4 的 Prop. 7 相同，只需考虑 §1, No.~8 的 Prop. 10。

现在转到 Prop. 13 的证明。

\begin{enumerate}
\def\labelenumi{\Alph{enumi})}
\tightlist
\item
  假设 \(\mathbf{X}\) 是紧的可度量空间，且 \(p\) 是连续满射：
\end{enumerate}

于是空间 T 是紧的可度量空间（GT, IX, §2, No.~10）。根据 Ch. VI, §3, No.~1 的 Th. 1，存在映射 \(t \mapsto \eta_t\)，从 T 到 \(\mathcal{M}_+(X)\)，它淡 \(\mu\) -可测并且标量意义本质 \(\mu\) -可积，使得 \(\nu = \int_{\mathrm{T}} \eta_t d\mu(t)\)，并且 \(\eta_t\) 的总质量为 1 且集中于 \(\overline{p}^{-1}(t)\)，对每个 \(t \in \mathrm{T}\) 均如此。令 \((\mathrm{S}_n)_{n \in \mathbf{N}}\) 是 T 关于 \(\mu\) 的一个分割，使得 H 在每个集合 \(\mathrm{S}_n\) 上的限制连续（§1, No.~8, Props. 10 and 11）；以 \(\Lambda: t \mapsto \lambda_t\) 表示从 T 到 \(\mathcal{M}_+(X)\) 的映射：它在 \(S = \bigcup_{n \in \mathbf{N}} S_n\) 上等于 H，在 T - S 上等于 0。显然，\(\nu = \int_{\mathrm{T}}\lambda_t d\mu (t)\)，且 \(\Lambda\) 满足叙述中的条件 \(a)\)。

设 \(\theta\) 是 T 上的测度；映射 \(\Lambda\) 淡 \(\theta\) -可测且标量意义本质 \(\theta\) -可积，因而也是 \(\theta\) -适配的（Ch. V, §3, No.~1, Prop. 2 b)）。令 \(f\) 是 X 上的普遍可测正函数；将 Ch. V, §3, No.~2 的 Prop. 5 应用于 \(\int \lambda_t d\theta(t)\)，可知映射 \(t \mapsto \lambda_t^{\bullet}(f)\) 是 \(\theta\) -可测的，从而由于 \(\theta\) 的任意性而普遍可测。

公式 (12) 由 Ch. V, §3, No.~2 的 Prop. 5 得到。

\begin{enumerate}
\def\labelenumi{\Alph{enumi})}
\setcounter{enumi}{1}
\tightlist
\item
  假设存在 \(\mathbf{X}'\)，它是 \(\mathbf{X}\) 的紧子集，集中着测度 \(\nu\)，且 \(p_{\mathbf{X}'}\) 连续：
\end{enumerate}

置 \(T' = p(X')\)，\(p' = p_{X'}\)；以 \(\nu'\) 表示测度 \(\nu_{X'}\)，并令 \(\mu'\) 为像测度 \(p'(\nu')\)，定义在 \(T'\) 上。由于 \(p'\) 是连续满射且 \(X'\) 是紧的可度量空间，根据 A)，存在映射 \(\Lambda': t' \mapsto \lambda_{t'}'\)，从 \(T'\) 到 \(\mathcal{M}_{+}(X')\)，满足以下条件：

\(a^{\prime})\) 对于每个 \(t^\prime \in \mathrm{T}^\prime\)，测度 \(\lambda_{t'}'\) 集中于 \(\mathrm{X}' \cap \overline{p}^{-1}(t')\)；

\(b')\) 对于每个普遍可测正函数 \(f'\)（定义在 \(X'\) 上），函数 \(t' \mapsto \lambda_{t'}^{\bullet}(f')\) 在 \(T'\) 上普遍可测，并且

\[
\int_ {\mathrm{X} ^ {\prime}} ^ {\bullet} f ^ {\prime} (x ^ {\prime}) d \nu^ {\prime} (x ^ {\prime}) = \int_ {\mathrm{T} ^ {\prime}} ^ {\bullet} d \mu^ {\prime} (t ^ {\prime}) \int_ {\mathrm{X} ^ {\prime}} ^ {\bullet} f ^ {\prime} (x ^ {\prime}) d \lambda_ {t ^ {\prime}} ^ {\prime} (x ^ {\prime}).
\]

令 \(t \in \mathbf{T}\)；如果 \(t\) 属于 \(\mathbf{T}'\)，则以 \(\lambda_t\) 表示 \(\lambda_t'\) 在 \(\mathbf{X}'\) 到 \(\mathbf{X}\) 的典范注入下的像；如果 \(t\) 属于 \(\mathbf{T} - \mathbf{T}'\)，则置 \(\lambda_t = 0\)。读者不难验证，映射 \(t \mapsto \lambda_t\) 满足叙述中的条件 \(a)\) 和 \(b)\)。

\begin{enumerate}
\def\labelenumi{\Alph{enumi})}
\setcounter{enumi}{2}
\tightlist
\item
  一般情形下的存在性：
\end{enumerate}

由于 X 上的测度 \(\nu\) 是适度的，我们可以取 X 的一个覆盖 \((\mathrm{U}_m)_{m\in \mathbf{N}}\)，它由 \(\nu\) -可积开集组成。此外，令 \((\mathrm{X}_n)_{n\in \mathbf{N}}\) 是 X 关于 \(\nu\) 的一个分割，使得 \(p\) 在每个集合 \(\mathrm{X}_n\) 上的限制连续（§1, No.~8, Props. 10 and 11）；以 \(\nu_n\) 表示 X 上的测度 \(\varphi_{\mathrm{X}_n}\cdot \nu\)，以 \(\mu_n\) 表示它在 \(p\) 下的像。根据 B)，对于每个整数 \(n\in \mathbf{N}\)，存在从 T 到 \(t\mapsto \alpha_t^n\) 的映射 \(\mathcal{M}_{+}(\mathrm{X})\)，满足以下条件：

\(a^{\prime \prime})\) 测度 \(\alpha_{t}^{n}\) 集中于 \(\overline{p}^{-1}(t)\)，对所有 \(t\in \mathbf{T}\) 成立。

\(b^{\prime\prime}\) 如果 f 是 X 上的普遍可测正函数，则 T 上的正函数 \(t \mapsto (\alpha_{t}^{n})^{\bullet}(f)\) 是普遍可测的，并且

\[
\int_ {\mathrm{X}} ^ {\bullet} f (x) d \nu_ {n} (x) = \int_ {\mathrm{T}} ^ {\bullet} d \mu_ {n} (t) \int_ {\mathrm{X}} ^ {\bullet} f (x) d \alpha_ {t} ^ {n} (x).\tag{13}
\]

有 \(\nu = \sum_{n\in \mathbf{N}}\nu_n\) 和 \(\mu = \sum_{n\in \mathbf{N}}\mu_n\)；由 No.~2 的 Prop. 3 及上述引理 4 立即可知，存在 T 上的普遍可测正函数序列 \((g_n)_{n\in \mathbf{N}}\)，使得 \(\mu_{n} = g_{n}\cdot \mu\) 对所有 \(n\in \mathbf{N}\) 成立，并且 \(\sum_{n\in \mathbf{N}}g_n = 1\)。对于每个 \(t\in \mathbf{T}\)，以 \(\beta_t^n\) 表示 X 上的测度 \(g_{n}(t)\cdot \alpha_{t}^{n}\)，以 \(q_{t}\) 表示 X 上的负荷 \(\sum_{n\in \mathbf{N}}(\beta_t^n)^{\bullet}\)。令 \(f\) 是 X 上的普遍可测正函数；利用 No.~2 的 Prop. 2 并对 (13) 中的 \(n\) 求和，得到

\[
\int_ {\mathrm{X}} ^ {\bullet} f (x) d \nu (x) = \int_ {\mathrm{T}} ^ {\bullet} q _ {t} (f) d \mu (t);\tag{14}
\]

此外显然，函数 \(t \mapsto q_t(f)\) 在 \(\mathbf{T}\) 上是普遍可测的。

对于每个 \(m \in \mathbf{N}\)，令 \(\mathbf{E}_m\) 为满足 \(t \in \mathbf{T}\) 且 \(q_t(\mathbf{U}_m) = +\infty\) 的点所成的集合；集合 \(\mathbf{E}_m\) 是普遍可测的，因为映射 \(t \mapsto q_t(\mathbf{U}_m)\) 具有此性质；并且 \(\mathbf{E}_m\) 对 \(\mu\) 局部可忽略，这是将公式 (14) 应用于 \(f = \varphi_{\mathbf{U}_m}\) 且 \(\nu^\bullet(\mathbf{U}_m)\) 有限的结果。于是集合 \(\mathbf{E} = \bigcup_{m \in \mathbf{N}} \mathbf{E}_m\) 普遍可测且对 \(\mu\) 局部可忽略。置 \(\lambda_t = 0\) 对于 \(t \in \mathbf{E}\)。此外，令 \(t \in \mathbf{T} - \mathbf{E}\)；负荷 \(q_t\) 由于开集 \(\mathbf{U}_m\) 覆盖 \(\mathbf{X}\) 且 \(q_t(\mathbf{U}_m)\) 有限而局部有界；根据 §1, No.~7 的 Prop. 7，存在测度 \(\lambda_t\)，定义在 \(\mathbf{X}\) 上，使得 \(q_t = \lambda_t^\bullet\) 且 \(\lambda_t = \sum_{n \in \mathbf{N}} \beta_t^n\)。映射 \(t \mapsto \lambda_t\) 满足叙述中的条件 \(a)\) 和 \(b)\) 是立即的。

D) c) 的证明：

设 \(f\) 是 \(X\) 上普遍可测的正有界函数；我们将证明，函数 \(h_f: t \mapsto \lambda_t^\bullet(f)\) 在 \(T\) 上普遍可测，并且是测度 \(\mu_f = p(f \cdot \nu)\) 关于 \(\mu = p(\nu)\) 的一个密度。令 \(K\) 是 \(T\) 的紧子集，并令 \(A = \overline{p}^{-1}(K)\)。对于每个 \(t \in T\)，测度 \(\lambda_t\) 集中于 \(\overline{p}^{-1}(t)\)；若 \(t\) 属于 \(K\)，则 \(\overline{p}^{-1}(t) \subset A\)，从而 \(\lambda_t^\bullet(f\varphi_A) = \lambda_t^\bullet(f)\)；另一方面，若 \(t\) 属于 \(T - K\)，则 \(\overline{p}^{-1}(t) \subset X - A\)，从而 \(\lambda_t^\bullet(f\varphi_A) = 0\)。将公式 (12) 应用于 \(f \cdot \varphi_A\)，\(^{(1)}\) 可得

\[
\mu_ {f} (\mathbf {K}) = \int_ {\mathrm{A}} ^ {\bullet} f (x) d \nu (x) = \int_ {\mathrm{K}} ^ {\bullet} d \mu (t) \int_ {\mathrm{X}} ^ {\bullet} f (x) d \lambda_ {t} (x) = \int_ {\mathrm{K}} ^ {\bullet} h _ {f} (t) d \mu (t),
\]

这就证明了关系 \(\mu_f = h_f \cdot \mu\)。

令 \(f = 1\)，可见函数 \(h_1: t \mapsto \| \lambda_t \|\) 是测度 \(\mu_1 = \mu\) 关于 \(\mu\) 的一个密度，因此在 T 中局部 \(\mu\) -几乎处处等于 1。

E) 唯一性：

令 \(t \mapsto \lambda_t^i\)（其中 \(i = 1,2\)）是从 \(T\) 到 \(\mathcal{M}_+(X)\) 的两个映射，满足叙述中的条件 \(a)\) 和 \(b)\)。如 C) 中一样，取 X 的一个 \(\mu\) -分割 \((X_n)_{n \in \mathbb{N}}\)，使得 \(X\) \(p_{X_n}\) 对每个 \(n \in \mathbb{N}\) 都是连续的，并置 \(N = X - \bigcup_{n \in \mathbb{N}} X_n\)。对于每个整数 \(n \in \mathbb{N}\)，取 X 上正函数的一个可数集合 \(D_n\)，这些函数在 \(X\) \(X_n\) 外为零，且它们在 \(X_n\) 上的限制在赋范空间 \(\mathcal{C}(X_n)\) 中构成稠密集（将 GT, \(X, \S 3\), No.~3 的 Th. 1 应用于可度量紧空间 \(X_n\)）。置 \(D = \bigcup_{n \in \mathbb{N}} D_n\)。

令 \(f \in \mathbf{D}\)；根据 D)，函数 \(t \mapsto (\lambda_t^1)^{\bullet}(f)\) 和 \(t \mapsto (\lambda_t^2)^{\bullet}(f)\) 是测度 \(\mu_f\) 关于 \(\mu\) 的密度，因此存在 T 中局部 \(\mu\) -可忽略集 \(\mathbf{E}_f\)，使得 \((\lambda_t^1)^{\bullet}(f) = (\lambda_t^2)^{\bullet}(f)\) 对于 \(t \in \mathbf{T} - \mathbf{E}_f\) 成立。此外，根据 (12)，令 \(\mathbf{F}_i\) 为 \(t \in \mathbf{T}\) 中满足 \((\lambda_t^i)^{\bullet}(\mathbf{N}) \neq 0\) 的点所成的集合；它是局部 \(\mu\) -可忽略的，对于 \(i = 1, 2\)。由于 D 是可数的，集合 \(G = \left( \bigcup_{f \in D} E_f \right) \cup F_1 \cup F_2\) 局部 \(\mu\) -可忽略；对于 \(t \in T - G\)，有 \((\lambda_t^1)^{\bullet}(\mathbf{N}) = (\lambda_t^2)^{\bullet}(\mathbf{N}) = 0\)，且

\((\lambda_t^1)_{\mathbf{X}_n} = (\lambda_t^2)_{\mathbf{X}_n}\)，从而根据 §1, No.~8 的 Prop. 9，有 \(\lambda_t^1 = \lambda_t^2\)。

证毕。

注记。--- 1) 如果 X 是苏斯林空间，则 X 的每个紧子空间都是苏斯林空间，因而是可度量的（TG, IX, Appendix I, Cor. 2 of Prop. 3），\(^{(2)}\)，并且 X 上的每个测度 \(\nu\) 都是适度的（§1, No.~9, Remark 1）。根据 Prop. 13，X 上的每个测度因此都相对于每个 \(\nu\) -逆紧映射具有一个分解。

\begin{enumerate}
\def\labelenumi{\arabic{enumi})}
\setcounter{enumi}{1}
\tightlist
\item
  沿用 Prop. 13 的记号，令 \(f\) 是正的 \(\nu\) -可测函数。可以如 Ch. V, §3, No.~2, Prop. 5 中那样证明：\(t \in \mathbf{T}\) 中使 \(f\) 不是 \(\lambda_t\) -可测的点所成的集合局部 \(\mu\) -可忽略，\(t \mapsto \lambda_t^\bullet(f)\) 是 \(\mu\) -可测的，并且关系 (12) 仍然成立。
\end{enumerate}

